\makeatletter
\def\input@path{{./}{../}{../../}{../../../}{../../../../}{preamble/}{../preamble/}{../../preamble/}{../../../preamble/}}
\makeatother

\input{preamble_exams.tex}
\usepackage{needspace}
\usepackage{titling}
\setlength{\droptitle}{-2em}
\pretitle{\par\begingroup\centering\LARGE}
\posttitle{\par\endgroup}
\preauthor{}\postauthor{}
\predate{}\postdate{}
\maketitlehookd{\vspace{-1em}}
\linespread{1.12}
\setlength{\parskip}{5pt}
\newgeometry{margin=0.65in,top=0.55in,bottom=0.65in}

\title{Grade 11 Term 2 Home Practice -- C4\\\large A mean with known variance\\[0.25em]
	\normalsize Global Economics -- 2026--2027}
\author{}
\date{}

\begin{document}
	\maketitle

	\setcounter{tocdepth}{1}
	\setcounter{secnumdepth}{2}
	\phantomsection
	\label{toc}
	\tableofcontents

	\section*{Evaluated criteria}
	\begin{enumerate}[label=\textbf{C\arabic*},start=4]
		\item Constructs confidence intervals for a population mean with known variance using the normal distribution, applying them to contexts such as average income, average investment return, or average price level.
	\end{enumerate}

	\textbf{How to practise.} Try each task before opening its blue solution link. The arrows return to the contents, problem list or current problem.

	\section{Worked Solution Examples}
	\label{sec:worked-solution-examples}

	\SectionProblemsTableOfContents{
		\SectionProblemLink{prob:c4-1}{Weekly earnings for courier shifts (C4)} \\
		\SectionProblemLink{prob:c4-2}{Warehouse handling prices (C4)}
	}

	% ----------------------------------------------------------------
	\subsection{Weekly earnings for courier shifts (C4)}
	\label{prob:c4-1}

	A courier cooperative draws a simple random sample of 64 distinct members from a roster of 2,400. Each member contributes their total earnings for the same reference week. The sample mean is 410 thousand COP. For this exercise take the population earnings variance to be known: $\sigma^2=6{,}400$ in squared thousand COP. Assume the distribution has no extreme skew or outliers, so the sample mean has an approximately normal sampling distribution at this sample size.

	\QuestionsTasks[sec:worked-solution-examples]{
		{4}/{sol:c4-1-a}/
		{Define the parameter, explain the normal procedure choice, and calculate the standard error from the supplied variance.},
		{4}/{sol:c4-1-b}/
		{Use $z^*=1.96$ to calculate and interpret the 95\% confidence interval. Would it support saying the population mean is above 400 thousand COP?}
	}

	\clearpage

	\Needspace{5\baselineskip}
	\phantomsection
	\label{sol:c4-1-a}
	\subsubsection*{C4 -- Select the procedure and compute the standard error}

	Let $\mu$ be the mean reference-week earnings of all 2,400 members on the roster. The population standard deviation is known: $\sigma=\sqrt{6{,}400}=80$ thousand COP. The random sample is under 10\% of the roster ($64/2{,}400\approx 2.7\%$), supporting approximate independence. The stated shape and sample size support the normal approximation for $\bar x$. Use the one-sample $z$ interval, with
	\[
	SE(\bar x)=\frac{\sigma}{\sqrt n}=\frac{80}{\sqrt{64}}=10\text{ thousand COP}.
	\]

	\secInicio[prob:c4-1][sec:worked-solution-examples]

	\Needspace{5\baselineskip}
	\phantomsection
	\label{sol:c4-1-b}
	\subsubsection*{C4 -- Construct and interpret a 95 percent interval}

	The margin of error is $E=1.96(10)=19.6$ thousand COP. Adding and subtracting it gives
	\[
	\bar x\pm E=410\pm19.6
	\quad\Longrightarrow\quad [390.4,429.6]\text{ thousand COP}.
	\]
	We are 95\% confident that the mean reference-week earnings for the roster fall in this range. The interval includes 400 and values below 400, so it does not establish that the population mean exceeds 400 thousand COP.

	\secInicio[prob:c4-1][sec:worked-solution-examples]

	\clearpage

	% ----------------------------------------------------------------
	\subsection{Warehouse handling prices (C4)}
	\label{prob:c4-2}

	\noindent
	\begin{minipage}[t]{0.53\textwidth}
	\vspace{0pt}
	An importer estimates the mean handling charge for pallets in a quarter. It draws a simple random sample without replacement of 36 receipts from a roster of 1,200, with one receipt per pallet. The sample mean is 74 USD per pallet. For this exercise use a known population variance of $144~\mathrm{USD}^2$ and assume the population of charges is approximately bell-shaped.
	\end{minipage}
	\hfill
	\begin{minipage}[t]{0.43\textwidth}
	\vspace{0pt}
	\centering\small
	Standard normal cutoffs, where the area shown is $P(Z\le z)$.\par\smallskip
	\begin{tabular}{@{}lccc@{}}
	\toprule
	Area to the left & 0.950 & 0.975 & 0.995 \\
	\midrule
	$z$ & 1.645 & 1.960 & 2.576 \\
	\bottomrule
	\end{tabular}
	\end{minipage}

	\QuestionsTasks[sec:worked-solution-examples]{
		{4}/{sol:c4-2-a}/
		{Construct a 99\% confidence interval for the mean charge. Use the table to find the cutoff.},
		{4}/{sol:c4-2-b}/
		{Explain what the interval estimates. Assess a proposed budget assumption that the quarter's population mean exceeds 80 USD per pallet.}
	}

	\Needspace{5\baselineskip}
	\phantomsection
	\label{sol:c4-2-a}
	\subsubsection*{C4 -- Select the critical value and construct the interval}

	The sample is random and $36/1{,}200=3\%$. The bell-shaped population supports a normal sampling model and $\sigma=\sqrt{144}=12$ USD is known, so use a one-sample $z$ interval for the population mean. As in Term 1, find the cutoff by area. A central 99\% leaves 0.005 in each tail, so the area to the left of the positive cutoff is 0.995 and $z^*=2.576$.
	\[
	SE=\frac{12}{\sqrt{36}}=2,\qquad E=2.576(2)=5.152,\qquad
	74\pm 5.152\ \Longrightarrow\ [68.85,79.15]\text{ USD (rounded)}.
	\]

	\secInicio[prob:c4-2][sec:worked-solution-examples]

	\Needspace{5\baselineskip}
	\phantomsection
	\label{sol:c4-2-b}
	\subsubsection*{C4 -- Interpret the endpoints for budgeting}

	We are 99\% confident that the mean handling charge across all 1,200 pallets in the quarter was about 68.85 to 79.15 USD per pallet. Every value in the interval is below 80 USD, so the sample gives evidence that the quarter's mean was below that benchmark. Individual pallet charges may exceed 80 USD.

	\secInicio[prob:c4-2][sec:worked-solution-examples]


\end{document}
